%! Parametric Functions Modeling Planar Motion — Unit 4, Lesson 4.2 — Warm-Up (Answer Key)
\documentclass[10pt]{article}
\usepackage{precalculus-article}
\usepackage{precalculus-key}   % pulls in -boxes; do NOT also load -boxes
\newcommand{\TallMath}[1]{$\displaystyle #1\rule[-1.4em]{0pt}{3.2em}$}

\begin{document}

\pageheader{Unit 4, Lesson 4.2}{Warm-Up --- Answer Key}
\namedateperiod

\vspace{0.10in}

\begin{notesbox}{Spiral Review --- Key}

\textbf{1.}\quad A particle's $x$-position is $x(t) = t^2 - 6t$.
Complete the table of values.

\vspace{6pt}
\begin{center}
\renewcommand{\arraystretch}{1.5}
\begin{tabular}{|c|c|c|c|c|c|c|c|}
\hline
$t$    & $0$ & $1$ & $2$ & $3$ & $4$ & $5$ & $6$ \\
\hline
$x(t)$ & \ans{$0$} & \ans{$-5$} & \ans{$-8$} & \ans{$-9$}
       & \ans{$-8$} & \ans{$-5$} & \ans{$0$} \\
\hline
\end{tabular}
\end{center}

\vspace{0.14in}
\textbf{2.}\quad Find the vertex of $y = -t^2 + 4t + 5$.
What is the maximum value of $y$?

\vspace{4pt}
Vertex at $t = -4/(2 \cdot (-1)) =$ \ans{$2$}.
Maximum value: $y(2) = -(4) + 8 + 5 =$ \ans{$9$}.

\begin{teachernote}
  Students should use the vertex formula $t = -b/(2a)$ explicitly.
  Reinforce that the maximum \emph{value} is $y(2) = 9$, not $t = 2$.
  This is the same technique they will use in the lesson to find
  vertical extrema of parametric motion.
\end{teachernote}

\vspace{0.10in}
\textbf{3.}\quad Solve $t^2 - 9 = 0$.

\vspace{4pt}
$t = $ \ans{$\pm 3$}

\vspace{6pt}
A particle moves with $x(t) = t^2 - 9$. What do the solutions above
tell you about the particle's position?

\vspace{4pt}
\ans{At $t = -3$ and $t = 3$, the particle has $x = 0$, so it is
on the $y$-axis at those moments. These are $y$-intercepts of the
parametric curve.}

\end{notesbox}

\end{document}
